% Standalone one-page pitch for ZQE applications
\documentclass[11pt]{article}
\usepackage[utf8]{inputenc}
\usepackage{amsmath,amssymb}
\usepackage{geometry}
\usepackage{hyperref}
\usepackage{parskip}
\geometry{margin=1in}
\title{ZQE: Applications and Practical Recipe}
\author{Will White}
\date{\today}

\begin{document}
\maketitle

\begin{abstract}
This note summarizes practical applications and an actionable recipe for ZQE
(``Zeroth‑order Quadratic Estimator'') as implemented in the project. The
document is a concise, self-contained briefing intended for a research
supervisor or collaborator: motivating idea, estimator definition, recommended
practical settings, and where ZQE is especially useful.
\end{abstract}

\section*{1. One‑Line Summary}
ZQE (``Consistent Decoder Estimation Without Accurate Posterior Inference'')
is a simple, simulation‑based objective that estimates decoder parameters
(loading matrix $W$ and intercept $b$) directly by matching decoder outputs
under real and simulated data using a fixed imputation of latent variables.
It enables consistent recovery of interpretable decoder structure even when
accurate posterior inference is difficult.

\section*{2. Formal Definition (short)}
Let $h_{\mathrm{enc}}(y)$ and $h_{\mathrm{dec}}(y)$ be chosen sufficient‑statistic
transforms of the observed response $y$. Let $\eta(z)=W z + b$ be the linear
predictor. The ZQE cross‑term objective optimises
\[
\mathcal{L}_{\mathrm{ZQE}}(W,b) \,=\,-\Biggl(\frac{1}{N}\sum_{i=1}^N
h_{\mathrm{dec}}(y_i)^\top\,\eta\bigl(\hat z(y_i)\bigr) 
- \frac{1}{M}\sum_{m=1}^M h_{\mathrm{dec}}(y^{\mathrm{sim}}_m)^\top\,\eta\bigl(\hat z(y^{\mathrm{sim}}_m)\bigr)\Biggr),
\]
where the imputed latent is the Gaussian MAP closed‑form:
\[
\hat z(y) = \bigl(W^{\top}W + \lambda I\bigr)^{-1} W^{\top}\bigl(h_{\mathrm{enc}}(y) - b\bigr).
\]
The outer optimisation updates $(W,b)$ (e.g. by SGD) while treating $\hat z(\cdot)$
as a fixed imputation (no gradient through the encoder). A centering term is
estimated by simulation from the generative model to remove nuisance terms.

\section*{3. Why this helps (intuition)}
- By fixing the imputation, gradients target only decoder parameters, avoiding
  the complexity of simultaneous posterior learning.
- The centering (simulation) term cancels biases due to the chosen
  $h$‑transforms and makes the objective robust across link choices.
- Closed‑form Gaussian MAP encoder is cheap and stable for linear decoders,
  and allows interpretable sparsity recovery in $W$.

\section*{4. Practical Recipe (recommended defaults)}
\begin{itemize}
  \item Transform choices: use $h(\cdot)=\log(1+y)$ (``log1p'') or the
    Anscombe/sqrt transforms for low counts; choose encoder/decoder transforms
    independently to match modelling assumptions.
  \item Warm start: initialise $W$ from the top‑$Q$ left singular vectors of
    centered $h(Y)$ (SVD warm start). This places parameters in a principal
    component rotation, accelerating sparsification.
  \item Proximal sparsity: apply an \emph{LR‑aware} proximal L1 operator
    (soft‑threshold after each optimiser step) with threshold $\tau = \lambda\cdot\mathrm{lr}$.
    This produces exact zeros, and scaling with the current learning rate makes
    sparsification aggressive early and fine later.
  \item Centering draws: draw a large number of simulated samples per outer
    update (we recommend at least $M\gtrsim 10N$ or a fixed minimum such as
    $M=2000$) under \texttt{torch.no\\_grad()} — this cheaply reduces centering
    variance without affecting parameter gradients.
  \item Numerical stability: add a small jitter to the encoder linear system,
    e.g. $W^{\top}W + (\lambda+\varepsilon) I$ with $\varepsilon\approx 10^{-6}$,
    and detect any fully‑zero loading columns after proximal steps and
    reinitialise them with tiny random noise to avoid singular solves.
\end{itemize}

\section*{5. Typical hyperparameters to try}
Start with these sensible defaults and then sweep if necessary:
\begin{itemize}
  \item Optimiser: SGD with an initial learning rate of $\approx 1.0$ and a
    Robbins–Monro decay schedule (flat‑phase then power decay).
  \item Proximal L1: $\lambda\in[0.001, 0.05]$ (larger values make sparser $W$).
  \item Centring draws: $M = \max(10N, 2000)$.
  \item Ridge (encoder): $\lambda_{\mathrm{ridge}}\approx 1.0$ plus jitter $10^{-6}$.
\end{itemize}

\section*{6. Applications}
ZQE is particularly useful where the decoder/observation model is the primary
object of inference and where posterior accuracy is costly or unnecessary:
\begin{itemize}
  \item Sparse factor loadings recovery (e.g. interpretable latent factors in
    single‑cell RNA‑seq or ecological trait models).
  \item High‑dimensional GLLVMs where posterior inference is expensive and a
    simple MAP imputation suffices to estimate a consistent decoder.
  \item Topic or latent factor models where direct decoder estimation (topics,
    item factors) is preferred to full posterior amortisation.
  \item Recommender or choice models where interpretable loadings are the
    primary target and exact posterior uncertainties are secondary.
\end{itemize}

\section*{7. Caveats and recommendations}
\begin{itemize}
  \item ZQE recovers the decoder under modelling assumptions on $h$ and the
    encoder; poor choices of transforms can slow or bias recovery — use simple
    transforms (log1p, sqrt) as baselines.
  \item If columns collapse to exact zero early (over‑aggressive prox), add
    small reinitialisation and reduce proximal strength or learning rate.
  \item Use SVD warm starts when seeking sparse solutions — random initialisation
    may require longer runs and more tuning.
\end{itemize}

\section*{8. How to compile}
On a machine with a standard LaTeX distribution (\texttt{pdflatex}):
\begin{verbatim}
pdflatex paper/pitch/zqe_applications.tex
\end{verbatim}
If you are on Windows and do not have LaTeX, installing MiKTeX is recommended.
On Linux, install TeX Live (e.g. \texttt{sudo apt install texlive-latex-recommended}).

\section*{9. Next steps / Offer}
If you would like, I can:
\begin{itemize}
  \item run a quick compile in this environment and return the PDF, or
  \item adapt the note to include a one‑page figure from your experiment
    (e.g. coefficient paths from a \texttt{lambda} sweep).
\end{itemize}

% Updated standalone one-page pitch for ZQE applications
\documentclass[11pt]{article}
\usepackage[utf8]{inputenc}
\usepackage{amsmath,amssymb}
\usepackage{geometry}
\usepackage{hyperref}
\usepackage{parskip}
\geometry{margin=1in}
	itle{ZQE: Applications and Practical Recipe (revised)}
\author{Will White}
\date{\today}

\begin{document}
\maketitle

\begin{abstract}
This one‑page brief summarises practical recommendations for the $Z_q$ (ZQE)
estimator based on the paper draft (`paper/main.tex`) and the project working
notes (`CLAUDE.md`). It emphasises the \emph{Gaussian proxy} as the canonical
practical choice, documents implementation safeguards observed in experiments,
and provides a concise recipe for reproducible experiments and sparse recovery.
\end{abstract}

\section*{1. Core message}
The $Z_q$ estimator yields consistent decoder estimation even with misspecified
or frozen surrogate inference, provided the centered estimating equations are
used and a local Jacobian nonsingularity condition holds (see `paper/main.tex`).
For practical use, the Gaussian proxy (choose $T(y)$ to make a homoskedastic
proxy) is the canonical instance: it is computationally trivial, theoretically
tractable, and empirically robust in large‑$p$ settings.

\section*{2. Precise estimator (practical form)}
Pick a transform $T(y)$ (e.g. $\log(1+y)$ or $\sqrt{y}$) and form the Gaussian
proxy so that
\[
T(y) \mid z \approx \mathcal N(Wz + b,\, \sigma^2 I).
\]
The closed‑form MAP encoder is
\[
\hat z(y) = (W^\top W + \sigma^2 I)^{-1} W^\top (T(y) - b),
\]
and the ZQE objective is the centered cross‑term
\[
\mathcal L(\theta) = -\Bigl(\mathbb E_{\text{data}}[T(Y)\cdot\eta(\hat z(Y))] - \mathbb E_{\theta}[T(Y^{\text{sim}})\cdot\eta(\hat z(Y^{\text{sim}}))]\Bigr).
\]
Compute the centering expectation by simulation (large fixed bank or on‑the‑fly
MC under \texttt{no\_grad()}). Do not backpropagate through the encoder solves.

\section*{3. Why the Gaussian proxy (summary of repo notes)}
- It gives a closed‑form encoder and often yields deterministic or low‑variance
  centering terms (multiple imputation with fixed seeds).
- In large‑$p$ regimes the ridge encoder concentrates (LLN in $p$), meaning the
  proxy approaches the posterior mean and ZQE recovers the complete‑data score.
- The Gaussian proxy lets us verify the Jacobian nonsingularity condition
  analytically in canonical settings (Poisson GLLVM), which strengthens the
  theoretical contribution for a statistical venue (JRSS‑B recommendation).

\section*{4. Implementation checklist (must‑do)}
\begin{itemize}
  \item \textbf{Encoder under no\_grad}: compute $\hat z(\cdot)$ inside
    	exttt{torch.no\_grad()} so gradients flow only through $\eta(\hat z)$.
  \item \textbf{SVD warm start}: initialise $W$ from top‑$Q$ PCs of centered
    $T(Y)$ to place parameters in a rotation where sparsity is discoverable.
  \item \textbf{LR‑aware proximal L1}: apply soft‑threshold after each
    optimiser step with threshold $\tau = \lambda \times \mathrm{lr\\_now}$
    (proximal ISTA scaled by current LR). This yields exact zeros and a
    meaningful regularisation path when warm‑starting large→small $\lambda$.
  \item \textbf{Large centering budget}: draw many simulated examples for the
    centering term (recommend $M=\max(10N,2000)$) inside \texttt{no\_grad()} to
    sharply reduce Monte‑Carlo noise at low cost.
  \item \textbf{Numerical safeguards}: add a tiny jitter $\varepsilon\approx 10^{-6}$
    when forming $W^\top W + (\sigma^2+\varepsilon)I$, and detect any fully‑zero
    columns after prox steps and reinitialise them with tiny Gaussian noise.
\end{itemize}

\section*{5. Practical defaults and tuning}
\begin{itemize}
  \item Optimiser: SGD with flat phase + Robbins–Monro decay; initial lr~1.0.
  \item Proximal L1: try $\lambda\in[0.001,0.05]$; warm‑start grid from large→small.
  \item Encoder ridge / proxy variance: $\sigma^2\approx 1.0$ (tune as needed).
  \item Centering draws: $M = \max(10N, 2000)$ as a robust default.
\end{itemize}

\section*{6. Experiments and story advice (paper/main.tex + CLAUDE.md)}
- Lead with the general ZQE framework and immediately present the Gaussian proxy
  as the canonical, constructive instance (answers the practical "what should I use?" question).
- Emphasise the large‑$p$ concentration phenomenon and report a p‑sweep figure.
- De‑emphasise the frozen‑amortised‑encoder experiment as a main claim (it
  creates interpretational complications); keep a short caveat in discussion.
- Include verification of the Jacobian nonsingularity for Poisson GLLVMs with
  the Gaussian proxy (this closes the main theoretical gap noted in project
  notes and helps target JRSS‑B).

\section*{7. Common pitfalls}
\begin{itemize}
  \item Do not backpropagate through the encoder; doing so reintroduces the
    chain‑rule terms the theory cancels and creates gradient leakage.
  \item L1 inside the loss (non‑proximal) rarely produces exact zeros — use a
    proximal soft‑threshold after the optimizer step to recover exact sparsity.
  \item If many columns die early, lower the prox strength or reinitialise
    dead columns — otherwise the encoder linear system becomes ill‑conditioned.
\end{itemize}

\section*{8. How to compile and share}
To compile locally with \texttt{pdflatex}:
\begin{verbatim}
pdflatex paper/pitch/zqe_applications.tex
\end{verbatim}
If compilation fails, install a TeX distribution (MiKTeX on Windows, TeX Live
on Linux). If you want, I can run a quick compile here and return the PDF.

\section*{9. Offer}
I can also:
\begin{itemize}
  \item insert a one‑page figure (lambda‑path coefficient heatmap or active‑count
    curve) into this note, or
  \item produce a short appendix summarising the Jacobian verification for the
    Poisson GLLVM with the Gaussian proxy (ready for JRSS‑B reviewers).
\end{itemize}

\vspace{1em}
\noindent Best,
\newline Will White

\end{document}
